\section{Data}
\label{sec:data}

\subsection{Source}

The target dataset is the Clothing Fit Dataset for Size Recommendation
(Kaggle: \texttt{rmisra/clothing-fit-dataset-for-size-recommendation})~\cite{misra2018decomposing},
covering garment reviews from online rental and retail marketplaces. Each
review carries static shopper body measurements (height, weight, bust size,
body type, age), an item identifier and category, the purchased garment
size, and the ordinal outcome \texttt{fit} $\in \{\text{small}, \text{fit},
\text{large}\}$.

\subsection{Synthetic fixture}

\texttt{scripts/download\_data.py} defaults to a synthetic generator (no
Kaggle account required, so the full pipeline -- including CI -- runs from a
clean checkout) that reproduces the dataset's structural properties rather
than its exact values:

\begin{itemize}
  \item \textbf{Clustered users and items.} Each user has a latent "true
        size"; each item has a nominal size and a fit sensitivity. Users
        self-select garments close to their own size (as real shoppers do),
        with enough spread that mismatches still occur regularly.
  \item \textbf{Skewed fit distribution.} The probability of \emph{small}
        or \emph{large} rises smoothly with the gap between a user's true
        size and the purchased item's nominal size, scaled by the item's
        sensitivity. The resulting label distribution is skewed toward
        \emph{fit} ($56.0\%$ fit, $23.3\%$ small, $20.7\%$ large in the
        fixture used for Table~\ref{tab:results}), matching the real
        dataset's class imbalance.
  \item \textbf{True chronology.} Every review carries a \texttt{review\_date},
        so the point-in-time machinery (Section~\ref{sec:feature-engineering})
        is exercised against real ordering, not an append-order proxy.
\end{itemize}

A \texttt{--kaggle} flag downloads the real dataset via the Kaggle API when
\texttt{KAGGLE\_USERNAME}/\texttt{KAGGLE\_KEY} are set in the environment;
credentials are never read from or written to the repository.

\subsection{Grouped user split}

Train, validation and test partitions are formed by randomly assigning
\emph{users} (not rows) to a partition, so every review by a given user lands
in exactly one partition
(\texttt{recommender.data.splitting.grouped\_user\_split}). This is distinct
from -- and in addition to -- the point-in-time cut-off described in
Section~\ref{sec:feature-engineering}: it specifically prevents a model from
learning "this exact combination of height/weight/body-type always runs
small" from a user's train-partition reviews and then trivially matching it
against the same user's test-partition reviews. \texttt{GroupSplit} raises if
construction ever produces an overlapping user across partitions, and this is
asserted directly in \texttt{tests/test\_features\_leakage.py}.
